% ssp245-review-summary.tex -- self-contained review summary for Charlie and Caspar
% Build: TEXINPUTS=/var/home/aoz/Dropbox/OZ-Labs/templates//: xelatex ssp245-review-summary.tex  (run twice)
\documentclass[11pt]{ozreport}

\renewcommand{\ozdocmark}{/var/home/aoz/Dropbox/OZ-Labs/templates/assets/oz-eye-gradient.pdf}

\title{Validation Review of the SSP2-4.5 Global\\ Temperature-Attributable Burden Run}
\author{Aaron Osgood-Zimmerman \quad\textperiodcentered\quad OZ Labs LLC}
\date{August 20, 2026}

\begin{document}
\maketitle

\begin{abstract}
We reviewed the first full production run of the global temperature-attributable
mortality pipeline (scenario SSP2-4.5, 204 countries, 27 climate models, years
2022--2050). All arithmetic and coverage checks pass. Against the published
2019 estimates of the methodology the pipeline implements (Burkart et al.,
\textit{The Lancet} 2021), our 2022 global total agrees to 0.1\%, country
ranks correlate at 0.993, and 159 of 204 countries fall inside the published
95\% uncertainty intervals. An apparent large shortfall in hot countries
relative to the newer GBD 2023 release turned out to be almost entirely
upward revision by GBD itself between the 2019 and 2023 releases, not a pipeline defect. The main
genuine limitations we identified and quantified are: risk curves that are
undefined above a daily mean of 32.1\,\textcelsius{} (a growing share of
future exposure), a cluster of tropical and small-island countries where our
burden runs low relative to GBD 2019 (temperature bias and damped daily
variability were tested and ruled out; the addendum identifies the
principal cause), and a years-of-life-lost definition that differs from
GBD's.
A post-review addendum (Section~\ref{sec:addendum}) documents a defect in
how uncertainty draws are paired, found after the main review; it warrants
rerunning ssp245 before launching the remaining three scenarios.
Recommendations are at the end.
\end{abstract}

\section{Background and scope}

The project estimates deaths and years of life lost (YLLs) attributable to
non-optimal ambient temperature, for 204 countries, projected to 2050 (and
later 2100) under emissions scenarios, for the World Bank. The methodology
follows Burkart et al.\ (2021), the study underlying the Global Burden of
Disease (GBD) treatment of temperature: for each of 17 temperature-related
causes of death, a set of exposure-response functions (ERFs; relative-risk
curves relating daily mean temperature to mortality) is combined with daily
temperature fields, a location-specific theoretical minimum risk exposure
level (TMREL; the death-weighted optimal temperature), and cause-specific
mortality, to produce population attributable fractions (PAFs; the share of a
cause's deaths attributable to non-optimal temperature) and attributable
deaths and YLLs. We use the ERF curves, TMRELs, and code released with the
2021 paper; mortality comes from Institute for Health Metrics and Evaluation
(IHME) national forecasts; daily temperature comes from the World Bank
Climate Change Knowledge Portal (CCKP) 0.25-degree daily product.

The run under review is the reference-scenario (SSP2-RCP4.5, ``ssp245'')
production pass executed by Caspar in July--August 2026: 204 countries
$\times$ 27 climate models $\times$ 29 years (2022--2050), 500 draws per
combination, 159{,}732 output files, summarized into review tables on
August 14, 2026. This document reports the review of that run, conducted
August 14--20, 2026. Unless stated otherwise, ``our'' numbers below are for
2022 and are built in two steps: each climate model's estimate is the mean
over its 500 draws, and we then average those 27 per-model estimates (an
unweighted ensemble mean).

\section{Data integrity}

The run's quality gates all pass: the additive identities (heat plus cold
equals non-optimal, for PAFs and deaths) hold to floating-point precision;
$|\mathrm{PAF}| \le 1$ everywhere (maximum 0.777); all values are finite; no
combination has an all-zero burden; coverage is complete (every country has
all 27 models and all 29 years; two further models publish no ssp245 data,
which is expected); and every one of the 159{,}732 combinations carries all
17 causes and YLLs.

\section{Validation against published estimates}

\subsection{The nine benchmark countries}

Table~1 of Burkart et al.\ (2021) publishes attributable deaths for nine
countries for 2019. Table~\ref{tab:nine} compares our 2022 ensemble against
it. Eight of nine of our point estimates fall inside the published 95\%
uncertainty intervals (UIs); the exception is Brazil (ratio 0.83), which is
low almost entirely on the heat side. The rank ordering across countries
matches (Spearman correlation 0.983). Note the comparison is not
year-matched: our 2022 (modeled climate, forecast mortality) against their
2019 (observed weather, GBD mortality); GBD's own numbers move little between
2019 and 2022, so the comparison is informative but should not be read to the
digit.

\begin{table}[b]
\centering
\small
\begin{tabular}{lrrr}
\hline
country & published 2019 [95\% UI] & ours 2022 & ratio \\
\hline
China          & 499{,}605 [425{,}094--573{,}429] & 555{,}481 & 1.11 \\
United States  & 110{,}588 [98{,}410--119{,}580]  & 103{,}964 & 0.94 \\
Mexico         & 19{,}055 [16{,}111--22{,}083]    & 20{,}787  & 1.09 \\
Brazil         & 17{,}286 [15{,}386--19{,}185]    & 14{,}383  & 0.83 \\
South Africa   & 8{,}815 [7{,}846--9{,}823]       & 9{,}575   & 1.09 \\
Colombia       & 4{,}820 [3{,}669--6{,}148]       & 4{,}746   & 0.98 \\
Chile          & 3{,}759 [3{,}244--4{,}267]       & 3{,}747   & 1.00 \\
Guatemala      & 1{,}258 [976--1{,}637]           & 1{,}095   & 0.87 \\
New Zealand    & 1{,}191 [996--1{,}361]           & 1{,}177   & 0.99 \\
\hline
\end{tabular}
\caption{Attributable non-optimal-temperature deaths per year: Burkart et
al.\ (2021) Table~1 values for 2019, with 95\% uncertainty intervals, against
our ssp245 ensemble means for 2022. ``Ratio'' is ours divided by published.}
\label{tab:nine}
\end{table}

\subsection{All 204 countries against the GBD 2019 temperature-attributable
burden estimates}

IHME published the GBD 2019 temperature-attributable burden estimates as a
dataset accompanying Burkart et al.\ (2021) (GHDx DOI 10.6069/7QWW-PF74):
attributable death rates and PAFs for all 204 countries for 2019. We
verified this dataset against Burkart Table~1 (PAFs agree to at most 0.015
percentage points, with identical UIs), reconstructed absolute counts from
its per-person rates (accurate to roughly $\pm$10\%, calibrated on the nine
Table~1 countries), and compared all 204 countries. Below, ``the GBD 2019
estimates'' refers to this dataset.

\begin{ozresult}{Headline validation result}
Against the GBD 2019 estimates, our 2022 global total agrees to 0.1\%
(1.684 vs.\ 1.686 million attributable deaths per year), the country
ranking correlates at 0.993 (Spearman), and 159 of 204 countries fall
inside the GBD 2019 95\% UIs. Counting a country as discrepant only when
the intervals on \emph{both} sides fail to overlap (their 95\% UI vs.\ our
95\% draw interval), a single country remains: Cape Verde. The pipeline
reproduces the methodology it implements at global and, for most countries,
national scale.
\end{ozresult}

Two disagreement counts are worth keeping separate. Our \emph{point
estimates} fall outside the GBD 2019 UIs in 45 countries; those decompose
as 41 low and 4 high, and of the lows 38 are tropical: a cluster of 18
small, mostly island countries (uniformly at 0.25--0.65 of the GBD 2019
value), and 23 larger countries (e.g.\ Indonesia at 0.56, Haiti at 0.35).
Section~\ref{sec:limits} reports the mechanisms tested. But once our own
uncertainty is included, only Cape Verde has non-overlapping intervals; the
other 44 are point-estimate differences within joint uncertainty. (Interval
overlap is a permissive test, since our draw intervals are wide where heat
and cold nearly cancel, which is why we report both counts.)

\subsection{GBD 2023, and an apparent hot-country shortfall}

Against the published 2022 estimates of GBD's newest release (GBD 2023),
our global total is 0.874 of theirs. The country-level agreement is
otherwise strong: ranks correlate at 0.991, and in 203 of 204 countries the
two estimates agree on the direction of the net burden (a country's net
non-optimal burden can be negative, i.e.\ net protective, where mild-day
protection outweighs harmful days; ours and GBD 2023's estimates put the
same countries on the same side of zero, with one exception). But there is
a marked pattern: the largest gaps are all hot, mostly lower-income
countries, with ratios of 0.38--0.58 (Niger, Sudan, Iraq, India, Nigeria,
Mali, Chad, among others).

We initially treated this as a possible pipeline defect and tested four
explanations (temperature bias, observed 2022 weather, mortality inputs,
risk-curve truncation). The decisive test became possible once the GBD 2019
estimates were in hand: the same countries compared against the \emph{GBD
2019} estimates are almost all inside those UIs (Niger 0.84, Sudan 1.31,
Kuwait 1.12, India 0.96, Iraq 0.94, Pakistan 0.90). GBD itself revised
these countries upward between the 2019 and 2023 releases by factors of
1.3--2.3 (e.g.\ Sudan 2.34, Niger 2.27, India 1.67).

\begin{ozresult}{Resolution of the hot-country gap}
The shortfall relative to GBD 2023 is, for these countries,
GBD's own between-round revision, not a defect in the run: the pipeline
matches the estimates of the methodology it implements. Which target the
project \emph{should} track (the implemented 2021 methodology, or the
GBD 2023) is a scoping decision, discussed in
Section~\ref{sec:rec}.
\end{ozresult}

Two country-level curiosities are worth recording. Cambodia appears 73\%
high against GBD's 2022 value but matches both GBD vintages' 2019 values;
GBD's own Cambodia estimate falls 42\% between 2019 and 2022 within the
GBD 2023 series, for reasons we have not identified. The United Arab
Emirates is the one hot country genuinely low against both GBD vintages:
our estimate is roughly half of the GBD value in each case (786 vs.\ 1{,}520
against GBD 2019 and 1{,}206 against GBD 2023). We ruled out temperature
bias there directly (its CCKP temperatures match the ERA5 reanalysis to
within 0.4\,\textcelsius{} in the mean), and its remaining gap is
unexplained at the point-estimate level, though our uncertainty interval
overlaps GBD's.

\section{Structural checks}

\textbf{Trends decompose sensibly once composition is removed.} Raw
attributable-death counts rise for both heat and cold in most countries
(global cold total 1.33M in 2022 to 1.99M in 2050). The cold rise looks
wrong at first: warming means fewer and milder cold days, so
cold-attributable deaths might be expected to fall. Two composition effects
account for it. First, population aging: the \emph{crude} cold share (cold
attributable deaths as a share of total deaths, with no adjustment for the
population's changing age mix) rises in 141 of 204 countries, but the
age-sex-\emph{standardized} share (the same quantity computed holding each
country's 2022 age-sex composition fixed) \emph{falls} in 140 of 204.
Second, the IHME mortality
forecast shifting deaths toward cold-heavy causes at high latitudes (chronic
kidney disease deaths in the high-latitude group grow from 269k to 687k by
2050). Within-cause, within-age climate signals behave as expected: the
age-standardized heat share rises in 197 of 204 countries. Any headline
trend statement should therefore be made on age-standardized attributable
\emph{fractions}, not raw counts.

\textbf{No climate model dominates the ensemble.} The models that sit
furthest from the ensemble median on average are the known
high-climate-sensitivity models (UKESM1-0-LL, HadGEM3-GC31-LL, KACE-1-0-G),
as expected. Strong single-model disagreements (one model far outside the
spread of the other 26) occur in 0.2\% of country-year combinations, and
inspection shows they concentrate where the burden is small (median case:
about 125 attributable deaths) or where the other models agree so tightly
that a modest deviation stands out; in the largest-burden such cases the
outlying model sits 10--20\% from the ensemble median (e.g.\ India 2034,
one model 20\% above), which the ensemble mean absorbs without difficulty.
A handful of near-zero-burden countries show
sign disagreement across models (e.g.\ Singapore, ensemble $-96$ deaths,
model range $-171$ to $+135$); for these, sign instability rather than
spread is the meaningful flag.

\textbf{YLLs per attributable death differ from GBD by construction, and
the life-table definition accounts for the size of the difference.} Our
global ratio is 13.9 years per attributable death versus 21.2 (GBD 2023)
and 19.4 (GBD 2019). GBD computes YLLs against its aspirational reference
life table (remaining life expectancy 88.87 years at birth, 21.29 years at
age 70), while we use each country's own United Nations World Population
Prospects life expectancy. To check that this alone explains the gap, we
applied GBD's reference table to our own attributable-death age
distribution: the result is 17.7--21.5 years per death (the range brackets
where deaths fall within our 5-year age groups), which spans GBD's
observed 21.2. It is a definitional difference, not an error, but any
side-by-side YLL comparison must state it.

\section{Quantified limitations}\label{sec:limits}

\subsection{Risk curves end at 32.1\,\textcelsius{}}

The Burkart ERF curves used in this analysis are estimated only over each
climate zone's observed daily-temperature range; the hottest zone's curve
ends at a daily mean of 32.1\,\textcelsius, and hotter days are evaluated
at the endpoint. IHME's released code does the same, so this truncation is
the published method. Its scale, measured from the CCKP daily fields
(two models): 7.9\% of world population lived in grid cells whose zone
assignment is clamped in 2022, rising to 14.9--23.5\% by 2050 under ssp245.
In the Gulf and Sahel, 25--40\% of person-days already exceed the curve
endpoint (Kuwait's exceedances average 5.5\,\textcelsius{} above it).

To bound what truncation could hide, we re-evaluated 2022 with the endpoint
log-linearly extrapolated (slope fitted to each curve's last
0.5\,\textcelsius): heat burden multiplies by 2.30 in Kuwait, 1.60 in
Niger, 1.31 in Sudan, and 1.00 in Haiti (which has no truncation). Under
that assumption, roughly half of the Gulf/Sahel gap to GBD
2023 (in log terms) closes.

\begin{ozcaveat}{Projection consequence}
Because the population share above the curve endpoints grows with warming,
projected burdens in hot countries are effectively lower bounds under the
published method: beyond the endpoint, additional warming adds no
additional risk. The extrapolation is an assumption, not evidence
(the true risk could plateau); we recommend reporting it as a labeled
sensitivity band alongside projection results, not adopting it silently.
\end{ozcaveat}

\subsection{The tropical low cluster and the exposure reference}

The temperature input deserves a precise description: the CCKP daily
product is NASA's NEX-GDDP-CMIP6, i.e.\ CMIP6 climate models bias-corrected
by quantile mapping to the GMFD observational dataset (Princeton's Global
Meteorological Forcing Dataset, reference period 1960--2014). It is neither
raw CMIP6 nor the ERA5 reanalysis that the 2021 estimates used for
exposure.

The countries where we run genuinely low against the GBD 2019 estimates
are overwhelmingly tropical. A natural suspect was damped day-to-day
temperature variability, since the burden lives in the spread of daily
temperature around the optimum. We tested it against ERA5 2022 in 17
countries, covering 8 of the 23 low-cluster members directly (Haiti,
Indonesia, Malaysia, Uganda, Tanzania, Honduras, Nicaragua, and the UAE).

\begin{ozcaveat}{The variance hypothesis did not survive the test}
Across the 17 compared countries, the CCKP-to-ERA5 ratio of within-pixel
daily standard deviation is 0.88--1.18 for 7 of the 8 low-cluster members
(Uganda, at 0.74, is the only marked deficit among them; Rwanda, not a
cluster member, is at 0.68), and the rank correlation between the spread
ratio and the burden ratio is near zero (0.11). Haiti
illustrates the point: spread ratio 0.92, burden ratio 0.35. Damped daily
variability is therefore not the general explanation for the tropical low
cluster. At the close of the main review the cluster's cause was
unidentified; the post-review audit subsequently traced its principal
cause to the uncertainty-draw pairing (Section~\ref{sec:addendum}).
\end{ozcaveat}

Two supporting observations. First, re-referencing the daily series to
ERA5 by per-pixel quantile mapping, tested on four countries, gave a mixed
verdict (it moves Rwanda onto the GBD 2019 value but moves the Philippines
\emph{away} from it), consistent with variance not being the driver; we do
not recommend it as a production change. Second, the exercise surfaced a
real sensitivity worth stating: the two observational references disagree
by 1.7\,\textcelsius{} in the Philippines' population-weighted mean
temperature, and switching between them moves that country's burden by
roughly a factor of two. No published uncertainty interval (ours or GBD's)
includes the choice of observational reference as a source of uncertainty.

\subsection{Other standing flags}

\emph{Sub-grid countries.} 29 countries are so small that at most one cell
center of the 0.25-degree temperature grid falls inside their borders, so
their entire climate signal comes from a single (usually ocean-dominated)
grid cell. Two symptoms follow: the 27 climate models disagree far more
about these countries than elsewhere (median across-model coefficient of
variation 0.50, versus 0.07 for resolved countries), and their estimated
attributable shares of death are much lower than comparable resolved
tropical countries (median 0.4\% of deaths versus 1.6\%). The cause of the
low shares is unresolved: the ocean-damping explanation was not supported
where it could be tested (on resolved countries, previous section), and
none of the 29 has an ERA5 extract for a direct test. Their numbers should
be flagged in any headline table. The 29:
American Samoa, Andorra, Antigua and Barbuda, Bahrain, Barbados, Bermuda,
Cook Islands, Dominica, Micronesia (Fed.\ States), Grenada, Guam, Maldives,
Malta, Marshall Islands, Monaco, Nauru, Niue, Northern Mariana Islands,
Palau, Saint Kitts and Nevis, Saint Lucia, Saint Vincent and the
Grenadines, San Marino, Seychelles, Singapore, Tokelau, Tonga, Tuvalu, and
the U.S.\ Virgin Islands.

\emph{Fixed TMREL.} The optimal temperature is held at its 2020 value for
all projection years (no adaptation). Defensible to 2050; a strong
assumption for 2100 that should be stated, with the pipeline's annual
TMREL derivation available as a sensitivity.

\emph{Forecast mortality.} Projected burdens mix the climate signal with
IHME's demographic and epidemiological forecasts; attributable
\emph{fractions} isolate the climate signal and should carry the headline
trends.

\section{Post-review addendum (August 20, 2026): uncertainty-draw pairing}
\label{sec:addendum}

After the main review, a line-by-line audit of the pipeline against IHME's
released code, prompted by the unexplained tropical cluster, found that the
cluster's principal cause is on our side, in how uncertainty draws are
paired.

The method propagates uncertainty as matched draws: for each draw, the
TMREL is defined as the temperature minimizing that same draw's
death-weighted risk curve; IHME's released TMREL-calculator code computes
exactly this. IHME's data release, however, contains 1{,}000
risk-curve draws but only 100 TMREL draws, with no index linking them, and
we verified the released TMREL draws do not correspond to the released
curve draws in any order. Our pipeline therefore paired them arbitrarily
(recycling the 100 across the 500 draws used in production). An
arbitrarily paired reference point is never the draw's own minimum, so
each mispaired draw's rescaled curve dips below 1 near its true minimum,
producing a one-signed downward bias, largest exactly where curves are
flat near the optimum: the tropical cluster's signature.

We confirmed this experimentally: a replica of the PAF stage using the
production pairing reproduces production essentially exactly (median ratio
1.00 across 158 countries; global total to 0.004\%), and switching it to
correctly derived per-draw TMRELs raises the low cluster substantially
(14 of 21 cluster members move inside the GBD 2019 UIs, from 5). The
correction also shifts overall levels up by roughly 7\%, leaving a
positive residual against GBD 2019 that we cannot yet attribute
(candidates include IHME's draw-specific TMREL weights, which are not
reconstructible from released artifacts, their whole-degree TMREL
rounding, and their subnational aggregation); agreement with GBD 2019
neither improves nor worsens on net, it redistributes. We judge the fix
warranted regardless: a bias with a proven sign and mechanism should not
be retained because it partially cancels against unidentified terms.

Separately, IHME's pipeline propagated daily temperature uncertainty
(their exposure data carries 1{,}000 draws built from ERA5's published
uncertainty field); the CCKP product provides no daily uncertainty, so our
run has no such term. Work is in progress to quantify what this omission
does to our results, using ERA5's measured uncertainty rather than a tuned
value, and to decide whether to add the term or to proceed without it as a
clearly labeled limitation. Three contained implementation bugs found by
the same audit (none affecting this report's numbers) have been fixed.

Consequence: the ssp245 burden step should be rerun with the pairing fix
and the temperature-uncertainty decision applied, before launching the
remaining scenarios. The temperature-conversion outputs are unaffected
and reusable.

\section{Recommendations}\label{sec:rec}

\begin{enumerate}
\item \textbf{Apply the draw-pairing fix and rerun ssp245's burden step
  before launching the remaining scenarios} (SSP3-7.0, SSP5-8.5,
  SSP1-2.6); see the addendum (Section~\ref{sec:addendum}). The
  temperature conversions are reusable, and the run machinery itself is
  validated, reproducible, and complete, so the rerun is a burden-step
  pass, not a restart.
\item \textbf{State the alignment target explicitly in all outputs}: these
  are estimates under the Burkart 2021 (GBD 2019-era) methodology. If
  alignment with GBD 2023 is later desired, the clean path is
  obtaining the GBD 2023 risk curves from IHME, not tuning ours.
\item \textbf{Report projected trends as age-standardized attributable
  fractions}, with counts secondary. This matches the original Workflow-B
  design, which built the deliverable around scenario comparisons and
  fractions precisely because calibration drift cancels there.
\item \textbf{Attach the truncation sensitivity to projection results} in
  hot countries (a two-model local computation suffices; no full-scale
  rerun needed).
\item \textbf{Flag the sub-grid 29 and the tropical low cluster} in
  country-level tables, with the exposure-reference caveat above.
\item \textbf{State the YLL definition} (national life expectancy, not the
  GBD reference life table) wherever YLLs are shown.
\end{enumerate}

\medskip
{\raggedright
\noindent\textit{Provenance: all numbers in this document come from the
review analyses in the project repository, run August 14--20, 2026,
applied to the production summary tables of August 14, 2026, the GBD
Results Tool exports of July 31, 2026, and the IHME GHDx dataset DOI
10.6069/7QWW-PF74. Repository files:}
\code{docs/reviews/ssp245-review-findings.org} \textit{and}
\code{output/review-ssp245/}\textit{.}
\par}

\end{document}
